% SPDX-FileCopyrightText: 2026 Bart van de Lint
% SPDX-License-Identifier: AGPL-3.0-or-later
\section*{Nomenclature}
\label{sec:nomenclature}

Terms of ship design and of airborne wind energy, in the sense this paper uses them. Each is also explained where it first appears.

\begingroup
\renewcommand{\theHtable}{nomenclature}
\footnotesize
\setlength{\tabcolsep}{4pt}
\begin{longtable}{@{}p{0.24\textwidth}p{0.72\textwidth}@{}}
\toprule
\multicolumn{2}{@{}l}{\textit{Ships and shipping}}\\
\midrule
Deadweight & The mass a ship can carry: cargo, fuel, stores and anything installed after delivery, such as a rig. A rig's mass therefore comes off the cargo of a ship that is full by mass. \\
Displacement & The total mass of the ship and everything on it, equal to the mass of water its hull displaces. \\
TEU, CEU & Twenty-foot equivalent unit, the space of one standard \SI{20}{ft} container; car equivalent unit, the deck space of one standard car. Capacities of container ships and car carriers are counted in them. \\
Slot, bay, tier & A slot is the space of one container; a bay is a cross-section of stacked slots across the ship; a tier is one layer of a stack. \\
Volume-limited ship & A ship that runs out of space before it reaches its deadweight, such as a container ship or a car carrier; a ship that reaches its deadweight first, such as a bulk carrier or tanker, is mass-limited. \\
Draught & The depth of the hull below the waterline. The \emph{design draught} is the loaded draught the engine and propeller are designed for; the deeper \emph{scantling draught} is the most the hull's structure allows, and sets the deadweight. In \emph{ballast} a ship sails without cargo and floats higher. \\
Air draft & The height of the ship above the waterline, which bridges and port cranes limit. \\
Block coefficient $C_B$ & How full the underwater hull is: its volume over that of the box of the same length, beam and draught. A tanker's is about 0.8, a fast container ship's about 0.65. \\
Prismatic coefficient $C_P$ & How the underwater volume is spread along the length: the volume over that of a prism with the largest cross-section, which is high for a hull with a long parallel middle. \\
Froude number Fn & The ship speed made non-dimensional, $V/\sqrt{gL}$; the waves a hull makes, and their resistance, scale with it. \\
Centre of buoyancy & The centre of the underwater volume; how far forward of the hull's middle it lies changes the wave resistance. \\
Leeway, heel & The small angle at which a hull moves sideways through the water to make the side force that balances a rig's side force; and the angle it leans over. \\
Service speed & The speed a ship is designed to make at sea with its engine at a stated part of its rating. \\
MCR & Maximum continuous rating, the most power an engine is rated to give continuously. \\
Sea margin & Extra power held back at the design point for the wind, waves and hull fouling of real service, as a share of the calm-water power. \\
Hotel load & Electrical power for the ship's own needs other than propulsion. \\
Diesel-electric drive & Several engines driving generators (gensets) that feed an electric propulsion motor, so that only as many run as the load needs. \\
VLSFO & Very low sulphur fuel oil, today's standard marine fuel. \\
e-ammonia & Ammonia made from hydrogen from renewable electricity and nitrogen from the air; a fuel without carbon, burned with a little fuel oil to ignite it. \\
e-methanol & Methanol made from hydrogen from renewable electricity and captured CO$_2$; a fuel without fossil carbon. \\
Class names & Size classes of trade: a \emph{coaster} sails short coastal routes; a \emph{feeder} brings containers to the large hubs; \emph{Panamax} and \emph{Neo-Panamax} ships are the largest that fit the old and the new Panama Canal locks; \emph{Megamax} the largest container ships; \emph{Handysize}, \emph{Ultramax}, \emph{Kamsarmax} and \emph{Capesize} bulk carriers run from about \num{40000} to \num{180000}~t; an \emph{MR} (medium range) tanker carries refined products, an \emph{Aframax} tanker and a \emph{VLCC} (very large crude carrier) crude oil; a \emph{PCTC} (pure car and truck carrier) carries cars. \\
Geared, gearless & A geared bulk carrier has its own cranes; a gearless one is loaded by the port's shiploaders. \\
Gross tonnage & A measure of a ship's enclosed volume, on which port dues are charged. \\
\midrule
\multicolumn{2}{@{}l}{\textit{Rigs and kites}}\\
\midrule
Rig & The wind propulsion system of a ship: a kite, a row of wing sails or a row of Flettner rotors. \\
Wing sail & A rigid upright wing standing on deck, trimmed to the wind like a sail. \\
Flettner rotor & A tall spinning cylinder on deck; the spin turns the wind into a large side force (the Magnus effect), at the cost of the power that spins it. The \emph{spin ratio} is its surface speed over the wind speed. In this paper a \emph{rotor} is always a Flettner rotor. \\
Apparent wind & The wind as felt on the moving ship or kite: the true wind less the velocity of the ship or kite. \\
Lift, drag, polar & The force of a wing across and along the wind it feels, each as a coefficient $C_L$, $C_D$ times the dynamic pressure and area; the polar is how they vary with the angle of attack, the angle between wing and wind. \emph{Stall} is the angle of attack above which lift falls. \\
Glide ratio & Lift over drag, $G=C_L/C_D$. A kite's pull grows with it, since it sets how fast the kite can fly across the wind. \\
Aspect ratio & Span squared over area; a slender wing has a high aspect ratio. \\
Projected area, flat area & The area of a curved kite seen from above, against the area of its canopy laid flat. All kite areas in this paper are projected. \\
Ram-air kite & A soft kite of two cloth skins with ribs between them, inflated by air entering openings at its front, as on a paraglider. Its fabric is the \emph{canopy}. \\
Tether, bridle, control unit & The single rope from the ship to the kite; the net of thinner lines that spreads its load over the canopy; and the pod at the end of the tether that steers the kite by pulling bridle lines. \\
Wind window & The quarter sphere downwind of the tether's anchor in which a kite can fly; at its \emph{edge} the kite hangs nearly at rest, in its centre it pulls hardest. \\
Figure eight, sweep & Flight patterns: a figure eight crosses the centre of the window repeatedly, which multiplies the kite's air speed and pull; a sweep moves slowly up and down along the window's edge. \\
Crosswind flight & Flight across the wind rather than hanging in it, which gives a kite an air speed several times the wind speed. \\
Depower & Lowering a kite's angle of attack so that it pulls less, to keep the tether force within its limit. \\
Launch, landing & Bringing the kite from its stowage on deck into the air, and back. \\
VTOL kite & A rigid wing on a tether with its own propellers, which takes off and lands vertically, hovering like a helicopter, and flies circles in between. \\
\midrule
\multicolumn{2}{@{}l}{\textit{Weather, routing and cost}}\\
\midrule
ERA5 & The European Centre for Medium-Range Weather Forecasts' reanalysis: past weather reconstructed on a global grid from observations and a weather model. \\
Significant wave height $H_s$ & The mean height of the highest third of the waves; the peak period $T_p$ is the period of the waves that carry the most energy. \\
Beaufort scale & A 0 to 12 scale of wind strength; Beaufort~3 is a gentle breeze of about 3.4 to \SI{5.4}{m/s}. \\
Weather routing & Choosing a ship's path and speed through forecast weather to lower its cost. \\
Cost per tonne-kilometre & The full cost of a round trip, capital, crew, port dues, fuel, kite wear and the interest on the cargo, over the cargo carried times the distance. \\
Saving & The drop in cost per tonne-kilometre that a rig gives against the same ship without a rig on the same fuel and routes, in per cent; differences between savings are in percentage points. \\
Goal & A value that development of a rig can still change, such as canopy life or crash rate, which the results assume is reached. \\
\midrule
\multicolumn{2}{@{}l}{\textit{Units}}\\
\midrule
kn, nmi & Knot, one nautical mile an hour, \SI{0.514}{m/s}; nautical mile, \SI{1852}{m}. \\
t & Tonne, \SI{1000}{kg}. \\
\bottomrule
\end{longtable}
\addtocounter{table}{-1}
\endgroup
